% ============================================================
% JMGM Conclusions — F3-b re-scope (2026-10-06); no numbers beyond Results.
% ============================================================

\section{Conclusions}\label{sec:conclusions}

Do semi-empirical NTO descriptors retain rank information on measured TADF
singlet--triplet gaps that structural fingerprints lose when molecules come from
publications absent from training? In this corpus they do: under source-paper
holdout the NTO model orders emitters across publications absent from training with
median Spearman $\rho = 0.34$, against $0.17$ for Morgan fingerprints. The
contribution is a matched three-unit holdout comparison showing that fingerprint
and RDKit ranking of measured TADF gaps depends on same-publication molecules in
training, while semi-empirical NTO ranking does not. Neither representation orders
molecules within a held-out paper, and neither has lower error than the other.
Three results delimit that answer.

\paragraph{Fingerprint and NTO models do not differ in error.}
Random forests built from the 2D structure alone---Morgan fingerprints
(MAE = \qty{0.091}{\electronvolt}) and RDKit descriptors
(\qty{0.100}{\electronvolt})---predict experimental $\Delta E_\text{ST}$ as well
as NTO descriptors derived from a semi-empirical excited-state calculation
(\qty{0.096}{\electronvolt}), against \qty{0.107}{\electronvolt} for the mean-value
predictor; the paired Morgan-minus-NTO MAE difference has a \qty{95}{\percent} CI of
\numrange{-0.0171}{+0.0092}~eV, and it stays near zero under paper holdout. Both
the Morgan and NTO models sit only \qtyrange{10}{15}{\percent} below the mean. This
null settles one design question: the sTDA calculation is not justified by
accuracy, and its measurable value is ranking across publications. The
\qtyrange{10}{15}{\percent} margin over the mean is the bar a new structure
representation has to clear on these labels.

\paragraph{The holdout unit, not the metric, sets the ranking contrast.}
The NTO-minus-Morgan $\Delta\rho$ is \num{0.16}
under source-paper holdout (\qtyrange{2.5}{97.5}{\percent} span \numrange{0.07}{0.25} over \num{50} fold
assignments) against \num{0.073} under scaffold and \num{0.055} under molecule
holdout. NTO ranks higher in \num{19} of \num{20} repeated paper partitions, but in
only \num{11} and \num{13} of \num{20} matched scaffold and molecule partitions, so
the \num{19}/\num{20} count belongs to the paper unit. The contrast comes from the
fingerprint side: NTO $\rho$ is nearly constant across units, while Morgan $\rho$
falls from \num{0.305} to \num{0.172}, and RDKit 2D descriptors fall with it
(\num{0.303} to \num{0.209}), so among the three sets tested the invariance is
specific to the NTO descriptors. Morgan ranks mainly molecules whose same-paper
siblings are in training; on molecules from single-molecule papers its $\rho$ is
\num{0.10}, and the NTO-minus-Morgan contrast on those molecules is not resolved
(\qty{95}{\percent} CI \numrange{-0.05}{0.37}). Siblings share both close structure
and measurement protocol, so this is an association, not an identified cause. SHAP
on the NTO variant ties the learned signal to hole--electron overlap and
charge-transfer character; it is an interpretable association, not a causal
validation.

Independent reports of the same compound
scatter by \qty{0.072}{\electronvolt} RMS. For the \num{102} multi-report molecules
the median label has an RMS standard error of \qty{0.049}{\electronvolt}; over all
\num{231} labels the estimated noise is about \qty{0.072}{\electronvolt} RMS, a
measurable part of the error budget but not all of it. Across \num{231}
molecules and \num{173} source papers, the report scatter supplies a direct
estimate of measurement variability. Error is lower on molecules whose
replicate reports agree (\qty{0.074}{\electronvolt} against
\qty{0.106}{\electronvolt}), but a constant predictor reproduces that ordering:
it reflects different spreads of the target, not evidence that removing label
noise alone improves the model. The estimated ceiling,
$R^2_\text{max} \approx 0.80$--\num{0.86}, lies well above the observed
$R^2 \approx 0.26$. In the tested representation set a frozen
pretrained molecular-transformer embedding reaches MAE
\qtyrange{0.109}{0.113}{\electronvolt}, worse than the Morgan model; other
representations remain untested. Separately, the computed gap is functional-dependent (mean
\qty{0.30}{\electronvolt}, up to \qty{0.58}{\electronvolt} between B3LYP and
CAM-B3LYP); this concerns computed references, not the experimental-label
noise estimate.

For modellers the consequence is practical. On literature-mined data, random and
scaffold splits leave same-paper siblings in training and, in this corpus, overstate
fingerprint ranking skill; a source-publication holdout should be reported with them
before a representation is chosen for ranking. Where ranking on unseen publications
matters, the semi-empirical NTO calculation is the representation that kept its rank
information here: it orders emitters across publications absent from training,
although no representation orders molecules within a held-out paper. As a triage filter the structure-based model enriches low-gap
candidates \numrange{1.2}{1.4}$\times$ over random selection at broad cutoffs; at
top-10 or top-20 the confidence intervals include random selection, and on the
filtered library the conformal intervals are too wide for molecule-level selection,
so no ranked shortlist is deposited.

All results come from one corpus of \num{231} molecules: the holdout-unit comparison
does not test transfer to a new laboratory or a controlled protocol, and
\num{128} of the \num{173} papers contribute a single molecule. The target is one
scalar, $\Delta E_\text{ST}$; the vertical NTO descriptors depart from adiabatic
references by \qty{0.195}{\electronvolt} on average for \num{14} emitters, and the
data do not support molecule-specific spin--orbit-coupling, $k_\text{RISC}$,
emission-wavelength or external-quantum-efficiency estimates
(\cref{subsec:limitations}).

\paragraph{What a controlled benchmark could resolve.}
Label noise of about \qty{0.072}{\electronvolt} RMS over the \num{231} molecules motivates
single-source, solvent-controlled measurements, with documented host and
protocol. Such a benchmark would separate measurement scatter from the
remaining modelling headroom; it does not follow from these data that a
better model could not improve on the current result.
